% CHR 2027 LIGHTNING TALK SUBMISSION
% Built from the official ACH LaTeX template (anthology-ch.cls, downloaded
% from https://anthology.ach.org/resources/ach-latex-template.zip).
% DO NOT change the documentclass. Add the 'final' option only for the
% camera-ready version -- the submission version renders "Anonymous
% Submission" automatically regardless of the \author fields below, per
% CHR's double-blind review policy.
%\documentclass[final]{anthology-ch} % for the final version
\documentclass{anthology-ch}         % for the submission

\usepackage{booktabs}
\usepackage{graphicx}

% TITLE OF THE SUBMISSION
\title{Training a Model to Disagree: Stream-of-Consciousness Typology as a
Site of Productive Ambiguity}

% AUTHOR AND AFFILIATION INFORMATION
% TODO: fill in real author name(s), ORCID(s), and affiliation(s) for the
% camera-ready version. Submission-mode rendering is anonymized automatically
% by the class file, so these values are safe to leave populated now.
\author[1]{Author One}[
  orcid=0000-0000-0000-0000
]

\affiliation{1}{TODO: Department, University, City, Country}

% KEYWORDS
\keywords[eng]{stream of consciousness, literary typology, large language
models, interpretive ambiguity, Virginia Woolf}

% METADATA FOR THE PUBLICATION -- leave defaults for submission
\pubyear{2027}
\pubvolume{1}
\pagestart{1}
\pageend{1}
\paperorder{1}
\conferencename{Proceedings of the Computational Humanities Research Conference 2027}
\conferenceeditors{}
\doi{00000/00000}
\customhead{}

\addbibresource{bibliography.bib}

\begin{document}

\maketitle

\begin{abstract}
Stream-of-consciousness (SoC) narration is a canonical case of a literary
technique that resists stable typology. We present preliminary results from
a pipeline that fine-tunes a local language model to detect and classify SoC
in Virginia Woolf's fiction against a taxonomy synthesizing Robert Humphrey
and Erwin Steinberg, and argue -- from a corpus-wide analysis of cross-model
disagreement -- that the persistent failure of frontier LLMs to agree on SoC
type is not noise to be filtered out, but a signal that tracks genuine,
scholarly-recognized ambiguity in the text. We describe a shift from
training on consensus to training a model that reports competing readings
directly, and situate this as a case study in modelling interpretive
ambiguity computationally.
\end{abstract}

\section{Motivation}

Stream-of-consciousness narration is a canonical test case for the limits of
literary typology: as Hoyt Long and Richard Jean So put it, ``no one can
agree on what it is, yet everyone seems to know it when they see it''
\cite{longso2016patternrecognition}. Decades of scholarship -- from Robert
Humphrey's foundational 1954 taxonomy onward -- have proposed competing
schemes for classifying SoC's techniques (interior monologue, free indirect
discourse, free association, orthographic markers, and more) without
settling where one shades into another \cite{humphrey1954stream}. We argue
that this persistent disagreement is not a defect to be engineered away, but
the most useful signal a computational method can capture about the form.

\section{Method}

This talk presents preliminary results from Penelope, a pipeline that
fine-tunes a local language model (Qwen3-4B) to detect and classify SoC
passages in literary prose against a ten-class taxonomy synthesizing
Humphrey and Erwin Steinberg \cite{humphrey1954stream, steinberg1973stream}.
Our corpus centers on Virginia Woolf's \textit{Mrs Dalloway} and \textit{To
the Lighthouse}, alongside a comparative set spanning Dorothy Richardson,
James Joyce, William Faulkner, and May Sinclair. Training labels come from
multiple frontier LLMs prompted independently on the same passages; a
consensus pipeline originally treated cross-model disagreement as noise to
be filtered out, keeping only passages where models converged on a single
type.

\section{Findings}

Corpus-level analysis of that discarded disagreement produced the finding
that reframed the project. Measuring cross-model agreement across fifteen
works, we found it never exceeds 27\%, and falls as low as 8\% for
Richardson's \textit{Pilgrimage} volumes. The overwhelming majority of
multi-model passage groups are \textit{disagreements}, not exceptions. When
we distinguish systematic disagreement (each model stably, repeatedly
assigning a different type across resamples) from random noise (a single
model flip-flopping between labels), a clear pattern emerges: the most
contested boundary across every author in the corpus is between
\textit{indirect interior monologue} and \textit{omniscient description} --
precisely the boundary literary critics identify as Woolf's most distinctive
and most contested technique, the point where free indirect discourse makes
it genuinely undecidable whether we are hearing the narrator or the
character \cite{zhang2019woolf}.

\section{Reframing the Model}

We treat this as evidence that model disagreement is tracking something real
about the text, not merely about model weakness. We are now training a
successor model (v3) to represent this directly: instead of collapsing a
passage to one label, the model learns to report \textit{both} defensible
readings, the textual evidence for each, and an explicit statement of where
and why the typology's boundary breaks down at that point in the passage.
This inverts the usual goal of a text classifier. Success is not measured by
whether the model converges with a single ground truth, but by whether its
stated uncertainty tracks the same passages where literary critics, and our
panel of frontier models, already disagree.

\section{Discussion}

We situate this as a case study in modelling interpretive ambiguity
computationally: rather than treating annotator or model disagreement as a
data-quality problem to be resolved by better prompting or larger consensus
panels, we ask what a classifier looks like when disagreement itself is the
object of study. We discuss the practical mechanics -- resampling frontier
models at temperature to distinguish stable from unstable disagreement,
representing dual readings in a training target, and reward-shaping a
reinforcement-learning stage to favor ``skeptical,'' evidence-grounded
reasoning over confident misclassification -- as well as the limits of this
approach, including the risk of a model that hedges indiscriminately rather
than locating ambiguity precisely. We close by asking what this suggests for
computational literary typology more broadly: if applying a systematic
typology to a whole text reliably produces ambiguity rather than resolution,
perhaps that instability is itself the most faithful thing a model can
report back.

\printbibliography

\end{document}
